%======================================================================
\subsection{BW-SIM-01 last step: thermal and thermal stress (plan, 30 Oct 2026)}\label{sec:sim-thermal-plan}
%======================================================================

The last two study types of \eref{SIM-01} reuse the frame that static, modal and buckling already use: the New
Study panel, Netgen mesh, bundled CalculiX, pre-check, results, bands and the \LaTeX{} report. The use case is
yours: an ESC or motor heating a printed mount, and whether the mount stays below the temperature where it
softens.

\paragraph{Theory.} Steady conduction solves $\nabla\cdot(k\nabla T)+\dot q=0$, with a fixed temperature,
a heat flux $q''$, or convection $-k\,\partial T/\partial n=h(T-T_\infty)$ on the boundary and the heat
generation $\dot q$ (W per body, spread over its volume) inside. CalculiX runs it as a \code{*HEAT TRANSFER,
STEADY STATE} step on the same C3D10 mesh, with $k$ from the Keystone cards. \textbf{Thermal stress} then
runs a static step with that temperature field as a load: each element expands by
$\boldsymbol\varepsilon_{th}=\alpha(T-T_{\text{ref}})\mathbf I$, so a free part grows and a held part is
stressed. It is one-way: the deformation does not change the heat flow, which holds for small strains.

\paragraph{Benchmarks (SIM-05).}
\begin{center}
\small
\begin{tabularx}{\linewidth}{@{}lX@{}}
\toprule
Bar, conduction & aluminium $10\times10\times100$\,mm, $k=167$\,W/m\,K, ends at 100 and 20\,$^\circ$C: linear,
  60\,$^\circ$C at mid-length, $Q=kA\Delta T/L=13.36$\,W.\\
Fin & same bar, base 80\,$^\circ$C, sides and tip in air at $h=25$\,W/m$^2$K, $T_\infty=20\,^\circ$C, tip
  insulated: $m=\sqrt{hP/kA}=7.738$\,m$^{-1}$, $T_{\text{tip}}=T_\infty+\theta_b/\cosh mL=65.64\,^\circ$C,
  $Q=\sqrt{hPkA}\,\theta_b\tanh mL=5.03$\,W; accepts 2\,\%.\\
Thermal stress & steel bar heated 50\,K: free, it grows $\alpha L\Delta T=0.060$\,mm; held at both ends,
  $\sigma=-E\alpha\Delta T=-126$\,MPa; accepts 2\,\%.\\
\bottomrule
\end{tabularx}
\end{center}

\paragraph{The workflow.} New Study offers \emph{Thermal} and \emph{Thermal stress}. The Setup group gains
Temperature, Heat flux, Heat source (W on a body) and Convection ($h$, $T_\infty$). The pre-check asks for $k$
(and $\alpha$ for thermal stress) and for at least one temperature or convection boundary, since without one
the temperature is undefined. The results show temperature in bands (like the safety factor) and, for thermal
stress, the static views. The report adds the heat balance: the heat put in equals the heat taken out by the
boundaries, to within the solver's tolerance.

\paragraph{Questions for you.}
\begin{enumerate}[label=H\arabic*.]
  \item \textbf{Study types now.} Steady-state thermal and one-way thermal stress, with transient (warm-up over
        time) later? \emph{Recommendation:} yes.
  \item \textbf{Loads.} Fixed temperature, heat flux, heat source in watts per body, convection; radiation later?
        \emph{Recommendation:} yes.
  \item \textbf{Convection presets.} Still air $h=10$, light airflow 25, fan or prop wash 75\,W/m$^2$K, plus a typed
        value? \emph{Recommendation:} yes, with a note that these are estimates.
  \item \textbf{Temperature limit.} Add each plastic's heat deflection temperature (HDT) to the Keystone cards, and
        colour the temperature result against it (amber within 10\,K of HDT, red above), like the safety-factor
        bands? \emph{Recommendation:} yes, configurable.
  \item \textbf{Parts in contact.} An ESC on a mount is two bodies. Treat touching faces as perfectly bonded for
        now (a contact conductance comes with \eref{SIM-02})? \emph{Recommendation:} yes.
\end{enumerate}
